\section{分母とその微分}\label{節:対数型の置換}

前節の公式で,特に $f(t)=\frac1t$ の場合を考えよう.一般的な形は
\[
 \int\frac{g'(x)}{g(x)}\dd x
\]
で与えられる.分子が分母の微分になっている分数である.\,$t=g(x)$ と置けば $\frac1t$ の積分になるので,対数に帰着させることを目標とする.

ここで1つ断りをいれておこう.本節の式は,\,$g(x)\ne0$ である区間で使う.\,$g$ が連続ならば,その区間内で符号を変えるには $0$ を通る必要があるから,符号は一定である.分母を $0$ にする点で区間が分かれ,任意定数は各区間で独立に選んでよい.

$\frac1t$ の原始関数は $\log|t|$ だから,\,$t=g(x)$,\,$\dd t=g'(x)\dd x$ と置いて
\begin{align*}
 \int\frac{g'(x)}{g(x)}\dd x
 &=\int\frac1t\dd t\\
 &=\log|t|+C\\
 &=\log|g(x)|+C
\end{align*}
を得る.実際,\,$g(x)>0$ の区間なら $\log g(x)$ を,\,$g(x)<0$ の区間なら $\log(-g(x))$ を微分すれば,どちらも同じ商 $\frac{g'(x)}{g(x)}$ に戻る.分子が $g'(x)$ の定数倍 $kg'(x)$ ならば,答えも $k\log|g(x)|+C$ となる.

\begin{bookpoint}
\textbf{POINT}\par
分子が分母 $g(x)$ の微分 $g'(x)$ の定数倍なら,\,$t=g(x)$ と置いて $\log|g(x)|$ の定数倍とする.
\end{bookpoint}

\begin{example}{絶対値を残すべき対数}
$\int\frac{6x}{3x^2-2}\dd x$ を求めよ.
\end{example}

\paragraphbf{方針}
分母 $3x^2-2$ の微分が分子の $6x$ なので,\,$t=3x^2-2$ と置いて対数にする.

\paragraphbf{解答}
$\log|3x^2-2|+C$

\paragraphbf{解法}
$t=3x^2-2$ と置けば $\dd t=6x\dd x$ なので,
\begin{align*}
 \int\frac{6x}{3x^2-2}\dd x
 &=\int\frac1t\dd t\\
 &=\log|t|+C\\
 &=\log|3x^2-2|+C
\end{align*}
となる.分母が $0$ になる $x=\pm\sqrt{\frac23}$ は除き,\,$x<-\sqrt{\frac23}$,\,$-\sqrt{\frac23}<x<\sqrt{\frac23}$,\,$\sqrt{\frac23}<x$ の各区間で考える.真ん中の区間では $3x^2-2<0$ なので,絶対値を外して $\log(3x^2-2)$ とすると実数の式ではなくなる.一方,\,$\log|3x^2-2|$ は3つの区間のどこでも使える形である.

\paragraphbf{部分分数との比較}
分母が二次式だから必ず部分分数にする,という決め方はしない.この例では分子を見れば分母の微分があり,1回の置換で終わる.前章の式の整理と置換積分は競合する規則ではなく,式の特徴に応じて選ぶ道具である.

ただし,\,$\int\frac1{3x^2-2}\dd x$ の分子は $6x$ の定数倍ではないから,同じ公式をそのまま使うことはできない.分母だけを $t$ と換えて,\,$\dd x$ を変えずに残すと,変数の異なる式を混ぜたことになる.置換後は,積分される式と微分の記号の両方を,新しい変数で書き直す必要がある.この積分は,分母を $(\sqrt3x-\sqrt2)(\sqrt3x+\sqrt2)$ と因数分解すれば,前章の部分分数で求められる.
